\section{Compositions discovered}
\label{sec:algos}

We use the word \emph{composition} rather than \emph{algorithm}
deliberately: AlphaEvolve, AlphaTensor, and FunSearch invent new
algorithms over open scoring functions, but the loop here does not
and cannot invent collective algorithms, because the underlying
\texttt{xm.*} primitives are black boxes whose schedules and chunking
are not reachable from above. The loop only chooses \emph{how many
primitive dispatches, in what order, against what tensor layout,
with what local pre- and post-processing}. What it emits for each
collective is one short Python file describing such a composition;
we summarize them below because the contrast with their
developer-written baselines is informative. Adding a further
collective is roughly $\sim$100 lines of glue: a developer
registers a new problem in \texttt{search/problems.py} with a
pure-Python reference, a small set of seed templates, and the I/O
shapes; Phases~1--5 then run unchanged. The cost model and the
profiling toolkit are shared across problems.

\paragraph{MoE token dispatch.} For \textbf{AllToAllV} the agent
packs the variable-length per-rank payload into a fixed-size buffer
where rank $i$'s payload starts at offset $i \cdot c_{\max}$, calls
one \texttt{all\_gather} to materialize a $(N, N \cdot c_{\max})$
tensor, and extracts each rank's receive slot with an
\texttt{index\_select} keyed on a trace-time-precomputed index vector
--- vs the two-dispatch AG$+$RS baseline Section~\ref{sec:agrs}
dissects. For \textbf{Uniform AllToAll} the agent emits a single
\texttt{all\_gather} followed by a \texttt{for}-loop over source
ranks that slices each destination's chunk with \texttt{narrow} and
concatenates via \texttt{torch.cat}; the contiguity-aware
implicit-copy term makes this win against the AG$+$transpose$+$RS
baseline, whose \texttt{permute$+$reshape} on a non-leading-dim
source pays a real $O(\text{numel})$ strided copy that the four-op
chain hides at the Python-op level.

\paragraph{Ring Attention KV.} For two-slot (K, V) distribution the
agent emits per-slot \texttt{all\_gather} calls --- two dispatches
per layer regardless of head count --- vs the developer-written
baseline that issues one \texttt{all\_gather} per head per slot
(16 dispatches per layer at HEADS=8, the OLMoE Ring KV probe shape).
Developers adopt the per-head form for principled reasons: small
collectives are typically easier to overlap with surrounding
compute, improve pipeline utilization, interact more predictably
with runtime fusion and network scheduling, and keep per-call
memory pressure low (the bundled form materializes a larger
contiguous intermediate, which under tight memory budgets risks
OOM). These are conservative defaults that generalize well across
cluster sizes and workloads.
The agent's contribution here is the empirical observation ---
surfaced by its profiled simulator --- that \emph{none of those
concerns bind} in this 7-node OLMoE setting: at the per-call
payloads here, network contention is unsaturated, HBM headroom is
ample, and the bundled intermediate fits comfortably. With those
constraints lifted, fewer-and-larger collectives reduce
per-dispatch launch and scheduling overhead. The per-call gain
is largest at 1-node bench scope, where the four-tier per-dispatch
launch floor dominates and Trainium has no opportunity to fuse
across the per-head calls: Table~\ref{tab:perproblem} reports
3.57\,ms baseline vs 1.07\,ms agent at 1-node bench
($3.34\times$). At 7-node bench scope the same per-call ratio
compresses to $1.46\times$ (2.95\,ms vs 2.02\,ms) because the
collective payload begins to share the per-call latency floor
with the cross-rank transport. At 7-node training scope the
per-step contribution is 13.49\,ms baseline vs 10.32\,ms agent
($1.31\times$), because the Neuron compiler's intra-graph
fusion within a single \texttt{mark\_step} already amortizes
part of the per-head dispatch floor that the bench scope cannot
share --- the agent's bundled form still wins, but the headroom
the developer baseline gives up at training scope is smaller than
at bench scope. The bundled form has a secondary advantage worth
flagging: it does not depend on the compiler's fusion behavior,
so it provides more stable per-call latency across vendor library
versions and avoids fragility when graph shapes shift across
runs.

\paragraph{Distributed cross-entropy.} For \textbf{distributed
cross-entropy} on vocab-sharded logits the agent computes
$\log\sum_v\exp(z_v)$ via two \texttt{all\_reduce}s. The first
is a \texttt{REDUCE\_MAX} over per-token local maxima, which
keeps the canonical max-shift required for bf16
numerical safety~\cite{kalamkar2019bf16, micikevicius2018mixed}.
The second is a single \texttt{REDUCE\_SUM} over a stacked
$(T, 2)$ tensor that packs both the shifted
$\sum_v\exp(z_v^{\text{local}} - \text{max})$ reduction and the
per-token in-shard target-logit reduction into one collective
dispatch. The non-obvious move is the stacking: a developer would
typically write three independent AR calls
(max, sum-exp, target); the agent collapses the two SUM ARs
into one by carrying them in the second dimension of a stacked
tensor.

The reference baseline this agent runtime is compared against in
Table~\ref{tab:perproblem} is the more common
distributed-CE realization: one \texttt{xm.all\_gather} of the
vocab-sharded logits along the head dimension, followed by a
single local \texttt{F.cross\_entropy}. A developer landing on
this 1-AG pattern is not being lazy. \texttt{F.cross\_entropy} does
the log-softmax, the max-shift, the padding-mask handling, and the
\texttt{ignore\_index} semantics in one
well-tested fused kernel with a matching backward;
a hand-rolled 2-AR composition has to re-derive each of those
correctly and trust bf16's exponent range, with the loss surface
silently breaking
(NaN, or worse a biased loss that still trains) if any of them is
wrong. The 1-AG pattern is also what NeuronX-Distributed and
Megatron-LM examples demonstrate, and it matches a single-rank
reference byte-for-byte, so it is the natural unit-test target
during convergence debugging. Notably, what the developer does \emph{not} get from picking 1-AG
is a per-call microbench advantage: the bench numbers in
Table~\ref{tab:perproblem} actually rank the deployed 2-AR agent
faster at both bench scopes ($4.5\times$ at 1-node, $3.3\times$
at 7-node), because the local $(T, V_\text{total})$
\texttt{F.cross\_entropy} the 1-AG path executes after the gather
is more expensive than the two scalar-tensor ARs of the 2-AR
path. The developer is trading per-call latency for correctness
guarantees, code idiom, and unit-test economy, and reasonably so
--- the per-call gap is a fraction of a millisecond, and the
correctness exposure of hand-rolling a stable distributed CE is
real.

The 2-AR agent runtime is also faster per-call than the 1-AG
developer-written baseline at both bench scopes
(Table~\ref{tab:perproblem}: 1.64\,ms baseline vs 0.37\,ms agent
at 1-node bench, $4.5\times$; 1.87\,ms vs 0.57\,ms at 7-node
bench, $3.3\times$). The win is volume: the two-AR composition
moves $(T, 2)$ scalars across ranks ($T = B \cdot S$) instead of
the $(T, V)$ tensor the 1-AG baseline materializes on every rank
before \texttt{F.cross\_entropy}, and at the small $T$ this
probe uses ($T{=}256$), one cross-rank move of $(T, 2)$ bf16
scalars is two orders of magnitude smaller than one
$(T, V{=}32256)$ bf16 tensor.

\paragraph{PP cross-stage.} The three Llama primitives of
Section~\ref{sec:eval-llama} share a structural template: $M$
independent per-microbatch collectives that the per-microbatch
\texttt{mark\_step} forces into $M$ separate HLO graphs, so the
Neuron runtime cannot fuse them. The baselines we compare against
are not strawmen but the developer realizations the AWS
NeuronX-Distributed tutorials demonstrate and that AWS
practitioners adopt --- those baselines are believed to be
best-performing on Trainium because they keep per-microbatch
activation working sets inside HBM, avoid compile-time graph
blow-up that NEFF compile latency would otherwise impose, and
rely on the in-graph fusion XLA delivers within each microbatch.
The agent's contribution is the negative observation that across
microbatches the per-graph dispatch floor compounds and a single
packed dispatch wins despite all three patterns. For PP
cross-stage the agent packs the $M$ per-microbatch payloads into
one leading-dim-extended tensor and issues a single dispatch:
a masked \texttt{xm.all\_reduce} over a
$(\textit{half}, M, \textit{Batch}, \textit{SeqLen}, D)$ buffer,
where \textit{half} is the rank count per pipeline stage
(here 112, since the 224 ranks are split into stage~0 and
stage~1) and the buffer simultaneously moves the $M$ forward
activations through the stage's send/recv pair. The agent
recognizes that masked-AR is the only correctness-preserving
substitute for \texttt{collective\_permute} on Neuron, where
the latter is broken. The structural cost is one large dispatch
instead of $M$ small dispatches, and the wall-clock saving comes
from collapsing $M$ Trainium per-dispatch latency floors into
one.

\paragraph{TP MLP.} The agent finds the same structural template
as for PP cross-stage but applied to the
$M \cdot N_\text{layers}$ partial-output tensors that the
per-microbatch \texttt{mark\_step} factors out across both
microbatches and TP layers. It flattens all of them into a single
contiguous 1-D buffer, emits one row-parallel AR, then
slices/reshapes back to the $[M][N_\text{layers}]$ structure the
caller expects. As with PP cross-stage the saving is one large
dispatch in place of $M \cdot N_\text{layers}$ small ones.

\paragraph{FSDP weight prefetch.} The interesting case is FSDP,
because the analogous ``stack-the-$M$-shards-and-issue-one-AG''
pattern that wins PP cross-stage and TP MLP \emph{does not
survive Phase~4's training-shape gate at 224 ranks}: the packed
shard tensor it constructs has shape
$(M \cdot N_\text{layers}, \text{ws}, \text{shard\_size})$ which
at the OLMoE Llama-block shapes exceeds the per-rank HBM
allocation budget the Neuron compiler tolerates when the
gathered intermediate is materialized in a single graph; the
attempted compile aborts before the per-step measurement can be
taken. This is the kind of hardware-binding constraint that
candidate-by-candidate static analysis cannot predict --- the
simulator's cost model ranks the packed candidate as the best,
and the Phase~3 LLM enumerates it confidently --- but the search
loop catches it at Phase~4 and falls back to the next-best
strategy among the candidates Phase~3 produced. That fallback,
which is what the deployed FSDP runtime actually is, is the
per-microbatch per-layer loop ($M \cdot N_\text{layers}$
individual AGs, structurally identical to the developer baseline
at the call level). The agent's contribution is therefore
\emph{not} dispatch collapsing; it is the choice to lift the
per-microbatch \texttt{mark\_step} that the developer baseline
emits between microbatches, letting NeuronCC fuse all
$M \cdot N_\text{layers}$ AGs into one HLO graph per training
step. The inline developer baseline keeps the per-microbatch
\texttt{mark\_step} for the reasons the introductory paragraph
of this section enumerates (HBM headroom, NEFF compile latency,
in-graph fusion within a microbatch), but at 7-node training
scale none of those concerns actually binds, and the agent's
single-graph variant collapses $M$ per-graph dispatch and
compile-amortization floors into one. The end-to-end ratio
(Table~\ref{tab:perproblem}, FSDP row: 8.52~ms baseline vs
2.01~ms agent at 7-node training) is the agent's contribution
even though no per-call collective count changes: the dispatch
count is identical, but the graph count drops from $M$ to $1$
per step. The broader point this FSDP case makes about the
search loop is the value of Phase~4 hardware validation:
the agent's enumerated theoretically-best candidate would have
crashed at compile if deployed blindly; the loop catches that
and converges instead to the next-best candidate that still wins
$4.2\times$ over the developer baseline. The agent's contribution
on a problem like FSDP is therefore as much the
\emph{robust convergence} to a deployable next-best as it is the
discovery of a structurally novel composition. The non-obvious common thread is that all three
problems look like fundamentally different parallelism patterns at
the framework level (pipeline, tensor, data) but reduce to the same
composition template once the per-microbatch \texttt{mark\_step}
structure is made explicit; the simulator's four-tier amortization
model (Section~\ref{sec:cost}) is what lets the agent see that
structure without running on hardware.

\FloatBarrier

